

\section{Sixth Sunday after Trinity: The Mind of the Flesh Is Enmity Against God}

\begin{quote}

Matthew 5:27–37. You have heard that it was said to those of old: You shall not commit adultery. But I say to you that everyone who looks at a woman in order to desire her has already committed adultery with her in his heart. But if your right eye causes you to stumble, tear it out and cast it from you. For it is better for you that one of your members perish than that your whole body be cast into hell. And if your right hand causes you to stumble, cut it off and cast it from you. For it is better for you that one of your members perish than that your whole body be cast into hell. 

It has also been said that whoever divorces his wife shall give her a certificate of divorce. But I say to you that whoever divorces his wife, except on account of fornication, causes her to commit adultery, and whoever marries a divorced woman commits adultery. Again you have heard that it was said to those of old: You shall not swear falsely, but shall perform to the Lord your oaths. But I say to you, do not swear at all: neither by heaven, for it is God’s throne; nor by the earth, for it is his footstool; nor by Jerusalem, for it is the city of the great King. Neither shall you swear by your head, for you cannot make one hair white or black. But let your speech be yes, yes; no, no. Whatever is more than this comes from evil.

\end{quote}

\bigskip


To admit that human nature is corrupted from birth, “conceived in sin and born in iniquity,” incapable of all good and inclined to all evil, will always be hard and distasteful to flesh and blood. For alongside the evil we think we see so much that is great and noble in the human spirit and in human life that the word, “that which is born of the flesh is flesh,” becomes to us altogether incomprehensible and unreasonable.

Must there really be no one among us who is good? Should it not be possible for someone, with the help of the good that is in him, if only he will in earnest, to combat and overcome the evil and thus become acceptable before God? Is not human nature in its essence good, because it is created by God, and has he not precisely given us his Law in order to show us how we may oppose and cast off all the uncleanness and sin that on account of our corruptibility presses in upon us?

These and similar questions have at all times been raised by all whenever the holy demands of God’s Law, with their sharp edge, have begun to stir the conscience. Then the excuses begin. To let the Law do its work, to acknowledge sin, and to come to God begging for grace—this the carnal ego resists and will rather resort to every invention and every illusion than to “give oneself up in order to gain oneself,” to “surrender one’s life in order to gain a life.”

And when then God’s Law stands just as holy, just as unyielding in its demands, and man neither dares reject it nor will bow beneath its judgment, then the question becomes how by every kind of trimming and reinterpretation one may break off its sharp edge and thus establish, as it were, a truce between the conscience and the Law, a righteousness of one’s own between God and man.

This is what we all, according to the carnality and dishonesty of our nature, are daily more or less consciously tempted to do, and it was precisely this in which the scribes of the Jews, after many hundreds of years of study of the Law, had become masters.

By all manner of interpretations, reductions, regulations, and traditions they had brought matters so far that they could live on excellent terms with the Law and, in their own opinion, fulfill all its demands and thereby gain a right to blessedness.

But the high thoughts of man must be brought low and his self‑deception uncovered if salvation is to come. Therefore the Savior strikes again and again with heavy blows of the two‑edged sword of the Law against the self‑righteousness of the Jews and rebukes their carnal and corrupt art of interpretation.

“How neatly you set aside God’s commandment in order to keep your own ordinance,” he says in one place (Mark 7:9), and here in our text he uncovers this hypocrisy of theirs in relation to the two things that most of all reveal the sinful and corrupted nature of man: lust and dishonesty.

At all times, and not least in our own day, men have labored hard to present lust or fornication as something purely natural, purely human, something that is by no means sinful in itself, even if for reasons of expediency it ought to be limited or regulated. This is a thoroughly heathen thought, and it is therefore no wonder that fornication in all its most abominable forms everywhere is the distinguishing mark and curse of heathenism.

But the Pharisees also, who boasted in the Law and acknowledged it as a command from God, “You shall not commit adultery,” sought to bring about a compromise between their lust and God’s holy will, partly by considering only the letter of the Law and partly by adapting its demands to the requirements of human and civil conditions, as it is also so beautifully said in our own day.

The Lord shows them that not only the outward act of adultery is sin, but even the “looking at a woman in order to desire her.” That is to say: our innermost thoughts, our most natural desires—lust itself as part of our being—is sin. “That which is born of the flesh is flesh.” 

How then can anyone stand? 

If the innermost thoughts and desires of my corrupt nature stand opposed to the very nature of the incorruptible God so that he must cast me away from himself in the judgment—where then is salvation? 

\textbf{At Golgotha there is salvation. God’s Word wounds, but it also heals.}

“For what the Law could not do, in that it was weak through the flesh, God did: sending his own Son in the likeness of sinful flesh and for sin, he condemned sin in the flesh, that the requirement of the Law might be fulfilled in us who walk not according to the flesh but according to the Spirit.”

God’s Word kills, but it also makes alive. And that which is born of the Spirit is spirit. He who in his distress has fled to Christ and through faith in him has received a new life, a new nature, does not steal away from the Law like the Pharisees and nominal Christians, but daily and in all earnestness allows the Law to condemn sin in his flesh, in order daily anew to flee to Golgotha and be washed clean in the blood of the Son of God.

Thus the Law on the one hand exposes sin, and on the other shows the believer what fruits faith must bear. This is what Jesus meant to teach the self-righteous Jews when he set forth the Law—and what his Holy Spirit teaches us daily.

But in this matter there is no compromise, no reconciliation with the world or with the so‑called purely human. It is precisely this purely human, this natural element, which in Scripture is called flesh, and of which it is said that the mind of the flesh is enmity against God.

Therefore Jesus rebukes the scribes for their careless and immoral view of the relation of marriage, namely that a man may divorce his wife if only he gives her a certificate of divorce. Assuredly the Pharisees maintained that it was necessary to adapt oneself thus far to the demands of imperfect human conditions and as it were regulate them and prevent greater excesses.

But that in its foundation this was only to reduce God’s demand in order to accommodate the lust of human nature and thus give ungodliness the appearance of righteousness—this the Lord clearly shows when he expressly declares that through such frivolous divorces there simply arises an adulterous relation.

This was surely hard for the Pharisees to hear and shameful, as when he asked that the one among them who was without sin should cast the first stone. It is therefore no wonder that in their bitterness they killed him; but his word was nevertheless truth.

One man and one woman. “The man shall leave his father and mother and cling to his wife, and the two shall be one flesh.” So it was from the beginning, and so it is God’s truth to this day, however much also the Pharisees and Sadducees of our own time seek to conceal it and by compromising with Satan and the world and the fornication of human nature dissolve the Word of God and even pronounce the church’s blessing over unions which the Lord himself calls adultery.

To divorce one’s wife is—like all other fornication—natural and human enough, but it is sinful and irreconcilable.

With the same uncompromising clarity the Savior speaks of the other evil mark of human nature: dishonesty. The lie is a part of our nature after him who is the father of lies. We lie to one another and to ourselves. Therefore men feel the need to swear, because they do not trust one another.

Should this be necessary among those who are born of God, who always remember his most holy presence? Should not their yes and their no be a holier assurance than the most solemn oaths of worldly men?

Therefore the Lord says: “You shall not swear at all”—a word which also in our time, when so many wish to mingle Christianity with authority and civilization, needs seriously to be laid to heart.

Cleanse yourselves. Flee from sin in every form. 

Flee from fornication. 

Become pure. 

Put away lying and speak truth with your neighbor, since we are members one of another. Add nothing to and take nothing from God’s holy Law lest he punish you. Do not be yoked together with unbelievers. For what fellowship has righteousness with unrighteousness? What communion has light with darkness?

Is not the mind of the flesh enmity against God?

